% !TeX root = ../papers/nc-p3.tex
\section{例外割線に接する平面局所点と三つの被約な割線}
\label{sec:finite-marked-tangent-secants}
\begin{theorem}\label{thm:finite-marked-tangent-secants}
定理\ref{thm:finite-marked-secant-pencil}で $\operatorname{length}\Sigma\le2$ なら、
印の接次元2でも全接方向の追加条件を除ける。
また $\Sigma$ が相異なる三点なら、接次元0または1の印で全く同じ反対対称性が成り立つ。
印以外の点はどちらの場合も非被約でよい。
\end{theorem}
\begin{proof}[平面局所環での二次の行関係]
長さ2で例外方向がある場合を扱う。印の局所環 $(S,\mathfrak m)$ の二つの生成元を $s,t$ とし、
例外接方向は $s=0$ とする。これは左右の各接写像が単射、例外割線の左右核が一意なので
ただ一つである。三座標を $l=s l_1+t l_0$, $r=s r_1+t r_0$ と書き、
係数ベクトル $l_i,r_i$ は $S$ 上に取る。前節の $A,B,Z=B-\lambda(B)A/\lambda(A)$ と
$N=\ker\lambda$ を使う。独立な形式変数 $v,w$ に関し、$N$ の行列を
$v l_1+w l_0$, $v r_1+w r_0$ で左右から縮約する六行三列を $H(v,w)$ とする。
実際の局所環では $H(s,t)Z=0$ である。

剰余体の射影直線上で $H$ の転置は $\mathcal O^6\to\mathcal O(1)^3$ を与える。
余核は $[0:1]$ にある長さ1の層であることを確かめる。
その点では階数2なので、局所的に余核は最大小行列式のイデアルによる環の商である。
そこで商が長さ2以上なら、その一次近傍の双対数環で $H$ の右核に、剰余が非零の $V\in N$ が取れる
（二つの主元で消去する）。$M=A\bmod\mathfrak m$ は全接方向の二つの縮約を消す。
従って $M,V$ が共通左右核を持つ相対鉛筆になり、前節の一意性からその核方向は一定である。
しかし $\bar l_0,\bar l_1$ は独立なので $v/w$ の一次変化で左方向は実際に動く。
これで長さ1が従う。核束は次数 $-2$、階数3で $\mathcal O^6$ の部分束なので、
分裂の各次数は $0,-1,-2$ のいずれかである。従って1だけひねった核の一次コホモロジーは消え、
像層の一次コホモロジーも $\mathcal O(1)^6$ の商であることから消える。
六行の一次式倍による三成分の二次式の余核は次元1で、例外方向の値で検出される。

$N^*$ の定数元 $e$ をこの方向の余核で非零に取り、上記二次式の余核を今度は $S$ 係数で作る。
剰余体上の結果と中山の補題で $w^2e$ がこの有限 $S$ 加群を生成する。
したがって、ある $a\in\mathfrak m$ について $vwe=a w^2 e$ が行を法として成り立つ
（例外方向で $v=0$）。各 $n\in N^*$ の $w^2 n$ も $c_n w^2 e$ に等しくなる。
基底に $e$ を含めて $c_e=1$ とできる。$(v,w)=(s,t)$ として $Z$ に作用させ、
$g=t^2 e(Z)$ と置くと
\[
 \mathfrak m^2\langle n(Z):n\in N^*\rangle=(g),\qquad
 t^2 Z=g V,\qquad (s-at)g=0.
\]
$V\in N\otimes S$ は基底に対する $c_n$ を成分に持ち、$e(V)=1$ である。
最後の式は $st e(Z)=a t^2e(Z)$ にさらに $t$ を掛けたものである。

$g$ 上では $l,r$ を $t(a l_1+l_0),t(a r_1+r_0)$ に置き換えられる。
$tg\ne0$ なら $\operatorname{ann}(tg)$ を法として、$A,V$ はこの二つの括弧に入る方向を
共通核に持つ。実際 $A$ の縮約は元から零、$V$ の縮約は $gV=t^2 Z$ から $g$ を掛けて零であり、
$t$ を取り出せばよい。両方向は剰余で非零、$A,V$ は剰余で独立である。
よって剰余長さにかかわらず例外割線の固定左核が得られ、$I_L g=0$ となる。
$tg=0$ なら $sg=0$ でもあり同じ結論となる。
前節の第二辺の評価 $f\in C$ は差 $Z$ の成分の局所線型結合なので、印では
$\mathfrak m^2 I_L f=0$ である。

印以外では既に $I_L f=0$。従って補題\ref{lem:finite-adjacent-word-support}の支持に
$L$ を課して失う長さは、印における $I_L f$ の長さに等しく、3以下である。
実際 $S$ は局所Gorensteinで $\operatorname{ann}(\mathfrak m^2)$ の長さは
$\operatorname{length}(S/\mathfrak m^2)=3$ である。
$f\ne0$ なら同補題により $b=3$ で、失う長さもちょうど3となる。
支持の長さは6以上ゆえ $Z(f)=\operatorname{Spec}(D/\operatorname{ann} f)$ の $L$ との交叉は
長さ3以上である。その第二射影は閉埋込みだから $L$ を張り、第二射影の包である平面は $L$ を含む。
この平面で $I_L$ は一本の線型式で生成され、特に印の $I_L f$ は単項である。
局所Gorenstein環の単項イデアルは零でなければその生成元の零化イデアルによる
Gorenstein商と加群同型である。ここではその極大イデアルの二乗が零なので、
socle次元1より長さは2以下である。これは必要な長さ3と矛盾する。
\end{proof}

\begin{proof}[三つの被約な階数一点]
三つの階数一行列の像ベクトルは独立である。
さもないと三行列はある共通左核を持ち、左縮約が階数1以下になる。
ここで三行列は独立である。三つが射影直線に入れば全二次小行列式がその直線で恒等零となり
$\Sigma$ の有限性に反する。右側でも同じ独立性が成り立つ。
従って三割線の左核は双対基底の三方向を与え、位置1では印を通る非共面な三直線 $L_1,L_2,L_3$ を得る。
局所Artin基底の相対鉛筆が階数一で切る次数2の有限平坦族は、三つの定数点の和のうち
二点の和の基底変換に限る。実際その三つの直和因子は各々平坦な環の直和因子で、
剰余で長さ1なら基底環と等しい。従って前節と同じ主元の消去、印での一様化変数の割算により、
評価 $f$ は各局所因子で三直線のどれか一つのイデアルで消される。点ごとに選ばれる直線は異なってよい。

二直線を含む平面の式 $h\in U_1$ を取ると、$hf$ の支持は残る直線に入る。
この支持に制限した第一座標の次元を $a'$ とすると、その長さは $a'$ 以下である
（補題の証明での $\mathbb P^{a'-1}\times\mathbb P^1$ 上の有限切断）。よって第一座標だけで
支持環へ全射する。一方 $\tau(hf U_0)=0$ は $f\in C$ の自己直交性から従う。
$hf$ が非零ならその支持商での完全跡に反する。三つの二直線の組の全てで $hf=0$ となる。
三平面の式は独立で共通零点は印だけなので、$f$ の第二射影の支持はスキームとして一つの被約点に含まれる。
補題の $b\ge3$ に反するため $f=0$ である。
\end{proof}
